% page 485 du fac-simile (image p-484 du scan)
% bloc 170.2 x 237.7 mm ; marge gauche 31.8 mm
\pgc{10.160mm}{0.000mm}{
\l{\cel{55}477}
\l{}
\l{\sou{querko}. (\sou{bot.})\cel{2}Granda arboro ek la klaso \textquotedbl{}amentacei\textquotedbl{}, di qua la}
\l{\cel{3}la ligno, tre harda, esas un ek le maxim uzata por menuzo e}
\l{\cel{3}karpento, e di qua la kortico, pulverigite, donas tano. -}
\l{\cel{3}L.\cel{1}\sou{quercus}. - L.}
\l{}
\l{\sou{questar}. (\sou{trans}.)Demandar, rekoliar, (pekunio) por verki karita-}
\l{\cel{3}tala o piesala. - FI.}
\l{}
\l{\sou{questionar}. (\sou{trans., pri}). Interpelar per demando di informivi}
\l{\cel{3}por igar su savoza pri ulo. - EFIS.}
\l{}
\l{\sou{questoro}. I. (en\cel{1}\sou{la epoki antiqua, en Roma)}.\cel{2}Oficiero qua ad-}
\l{\cel{3}ministris financi publika. - II. (\sou{nun}).La persono qua, en ula}
\l{\cel{3}korporacioni, esas komisita por surveyar la uzo di la pekunio}
\l{\cel{3}korporacionala. - DEFIRS.}
\l{}
\l{\sou{qui}.\cel{2}(\sou{gram.})\cel{2}Pluralo di \textquotedbl{}qua\textquotedbl{}.}
\l{}
\l{\sou{quieta}. Di qua la psiko ne agitesas da angori. - EFIS.}
\l{}
\l{\sou{quik}. Sen ajorno, tardeso (+retardeso) o fristo, mem minima.}
\l{\cel{3}E.}
\l{}
\l{\sou{quilayo}. (\sou{bot.}) Genero de rozacei, ek tri o quar speci de arbori,}
\l{\cel{3}kun folii alternanta, flori dioika, qui kreskas en la sudo-}
\l{\cel{3}regioni di Amerika, e di qui la pulvero, obtenata de lia kor-}
\l{\cel{3}tico, uzesas por netigar la linji sordida de makuli grasa.}
\l{\cel{3}- L.\cel{1}\sou{quillaja saponaria}. -}
\l{}
\l{\sou{quino}. (\sou{bot.)(farmacio}). Kortico kontre-febra, donata da la ar-}
\l{\cel{3}bori exotika ek la genero \textquotedbl{}cinchona\textquotedbl{}. DEFIRS.}
\l{}
\l{\sou{quin}.\cel{2}Pluralo di \textquotedbl{}quan\textquotedbl{}.}
\l{}
\l{\sou{quinara}. (\sou{matem.})\cel{2}En qua kin unaji simpla formacas unajo ek}
\l{\cel{3}la klaso kinesma. - DEFIS.}
\l{}
\l{\sou{quingo}. (\sou{bot.)}\cel{2}Frukto flava, di qua la pulpo kontenas quaza}
\l{\cel{3}stoneti, acerba ed astriktiva, - donata da arboro ek la fami-}
\l{\cel{3}lio \textquotedbl{}rozacei\textquotedbl{}, L.\cel{1}\sou{cydonia}, kun folii di qui la suba facio aspek-}
\l{\cel{3}tas kotonoza.}
\l{}
\l{\sou{quinino}. (\sou{kemio)}. Alkaloido de vejetanto, extraktita de quino.-}
\l{\cel{3}\sou{Simbolo kemiala}\cel{1}: C20 H24 O2 N2. DEFIRS.}
\l{}
\l{\sou{quinkunco}. Asemblitaro de objekti (ordinare : arbori) dispozita}
\l{\cel{3}kinope : quar qui formacas la anguli di quadrato, ed un en}
\l{\cel{3}la centro. - DEFI.}
\l{}
\l{\sou{quinqua}-.. .\cel{2}Prefixo ciencala : qua havas kin...}
\l{}
\l{quinquagezimo. (\sou{liturgio katol.}) La kinadekesma dio ante la}
\l{\cel{3}pasko-festo. - DEFIS.}
}
